\subsection{A third scalar input and batched preprocessing}
\label{sec:rank-three-critical}

The preceding analytic argument extends to every fixed first-layer
rank. Preprocessing is the restriction: the known function must be
compiled on a multidimensional grid. Batching all grid centers before
each dense linear layer brings rank three within the original budget.

\begin{theorem}[Critical sampling after a rank-three first layer]
\label{thm:rank-three-critical}
Under the assumptions of Theorem~\ref{thm:rank-two-critical},
replace the first-layer rank bound by three. Exact sampling uses
$O(\sqrt{D+1})$ expected input probes and
$O(\sqrt{D+1}\,B^5)$ expected online bit operations.
At $D=\lceil\log_2n\rceil$ and $b=O(\log n)$, preprocessing
and storage remain $O(Dn^2\log^6n)$. The worst-case query order
is $\Theta(\sqrt D)$ under the same width and grid conditions.
For first-layer rank at most two, the sharper $B^4$ online bound
of Theorem~\ref{thm:rank-two-critical} still applies.
\end{theorem}

\paragraph{Proof idea.}
The analytic estimates and source law extend to three variables. Their
direct dense preprocessing would exceed the chosen logarithmic budget. Each
linear layer, however, acts on the Taylor coefficients at every grid center
by the same weight matrix. Compute all those products together using exact
integer matrix multiplication. The seven-product recursion supplies enough
saving even with schoolbook integer arithmetic.

\begin{proof}
We give the construction for any fixed $r$, then specialize its
preprocessing count to $r=3$. The determinant-doubling basis
construction gives $r$ original rows of absolute sum at most one
and representation coefficients of magnitude at most two. Its
fixed-order rational arithmetic has $O_r(B)$-bit operands and
$O_r(n^2B^3)$ work. Thus $r$ independent scalar source streams
represent the first layer, with at most one input probe per draw.
The first preactivation on the enlarged box is bounded by $c=2r$.
After the first tanh, every output has magnitude at most one, so
the subsequent critical radius bound is unchanged.

For $D\ge2$, use a dyadic gate $q$ with
$q\le6/\sqrt{D+1}$, $q^{-1}\le\sqrt D$, and
$\|F/q\|_\infty\le4/5$. Depth one has the direct unit-row
factory. The derivative induction in the rank-two proof holds at
every order with first-preactivation bound $c=2r$. With the same
constants $C_k$, it gives $\|D^kF\|_{1,\mathrm{ent}}
\le C_k(2r)^kD^{(k-1)/2}$ through order $2r+2$.
For $f(p)=F(2p-1)/q$, take
\[
 M_k=2^kC_k(2r)^k(D+1)^{\lceil k/2\rceil}.
\]
In Lemma~\ref{lem:fixed-rank-bernstein},
$S_\ell=O_r(D^{\ell+2})$, $m_0=O_r(D+1)$, and
$\overline A=O_r((D+1)^2)$. The gate therefore gives
$O_r(\sqrt{D+1})$ expected probes.

The uniform Taylor oracle uses a complex tube of width
$1/(8r\sqrt D)$, grid spacing $\Theta_r(D^{-1/2})$, and
$O_r(D^{r/2})$ centers. A radius-two polydisk in scaled variables
lies inside this tube. Each output coefficient satisfies
$|b_\nu|\le2^{1-|\nu|}$. A total-degree-$K$ truncation, with
$K=O_r(B)$, provides all required derivatives through order $2r$.
Its geometric tails have a fixed polynomial factor in $K$; the
only scale conversion is at most $a^{-2r}q^{-1}$, polynomial in
$D$ for fixed $r$. Choose
$P=\lceil2\log_2(n^2D)\rceil=O(B)$ and
$\varepsilon=2^{-P}\le(n^2D)^{-2}$.
Thus $O_r(B)$ stored bits per coefficient suffice for this
absolute accuracy $\varepsilon$.
The total-degree residual calculation also has explicit finite
precision margins. On the unit polydisk, neuron coefficient norms
are at most $2^{r+1}$. First-layer preactivations are affine with
coefficient norm at most $2r+1$, and later preactivations have norm
at most $2^{r+1}$. Thus a perturbed exponent input has norm at most
$2^{r+2}+1$. Every complex preactivation has imaginary part below
one half, so $|1+e^{2s}|\ge1$ and reciprocal coefficient norm is
at most $2^r$. Put
\[
 a_r=2^{r+4}+4r+16.
\]
The exponential residual identity from the rank-two proof bounds
nonlinear-layer amplification by $2^{a_r}$ and local rounding error
by $2^{2a_r}n(K+2)^{4r+8}2^{-M}$. Indeed the Euler recurrence has
at most $(K+1)^{2r}$ products; its degree factors and divisions
contribute only polynomial powers of $K$. Solving its residual
equation has norm at most $\exp(2(2^{r+2}+1))$. An inverse
perturbation of coefficient norm below $2^{-r-1}$ keeps the inverse
norm at most $2^{r+1}$. These bounds fit the displayed constants.
If the stored coefficient precision is $S_0$, it suffices to take
\[
 M=S_0+4a_r(D+1)+
 \lceil\log_2(256nD(K+2)^{4r+10})\rceil+4rb.
\]
The resulting error induction verifies all perturbation margins
and leaves width below $2^{-S_0}$. Here
$S_0=O_r(B)$ and hence $M=O_r(B+D)$. There are
\[
 N_K=\binom{K+r}{r}=O_r(K^r)
\]
coefficients per center. Rounded evaluation of the required fixed
list of derivatives costs $O_r((B+\operatorname{bits}(p))^{r+2})$.

Each factory correction contains only one $O(m)$ hypergeometric
average. The rank-two uniform-comparison argument, with the evaluation
power $r+2$ in place of four, therefore gives
$O_r((D+1)B^{\max\{4,r+2\}})$ conditional expected work, including source
selection, coefficient evaluation, and arbitrary precision
refinement. For $r\ge2$, the direct augmented-forward cost with its fourth
power of precision is included in the exponent $r+2$.
Multiplying by the gate gives the stated online bound for $r=3$.

It remains to charge preprocessing. Batch every Taylor coefficient
at every center into one matrix $X$ with $n$ rows and
\[
 C=O_r(D^{r/2}K^r)
\]
columns. At each subsequent linear layer compute $WX$ together.
The entries of $X$ have a common working dyadic denominator $2^M$,
and the original weights have denominator $2^b$. Clear these
denominators, multiply the resulting integer matrices exactly, and
round the completed entries back to the working grid. There is no
rounding inside the matrix multiplication.

We use Strassen's seven-product recursion~\cite{strassen1969}.
For clarity, its exact arithmetic identities are recorded here.
For two block matrices $A,B$, form
\begin{align*}
 Q_1&=(A_{11}+A_{22})(B_{11}+B_{22}),&
 Q_2&=(A_{21}+A_{22})B_{11},\\
 Q_3&=A_{11}(B_{12}-B_{22}),&
 Q_4&=A_{22}(B_{21}-B_{11}),\\
 Q_5&=(A_{11}+A_{12})B_{22},&
 Q_6&=(A_{21}-A_{11})(B_{11}+B_{12}),\\
 Q_7&=(A_{12}-A_{22})(B_{21}+B_{22}).
 \end{align*}
The output blocks are
\[
 C_{11}=Q_1+Q_4-Q_5+Q_7,\quad C_{12}=Q_3+Q_5,\quad
 C_{21}=Q_2+Q_4,\quad C_{22}=Q_1-Q_2+Q_3+Q_6.
\]
Direct expansion proves these identities over the integers. The
recurrence $T(k)=7T(k/2)+O(k^2)$ gives $O(k^\omega)$ scalar
operations, $\omega=\log_2 7$.

Choose a power of two $k\in[C,2C)$, pad $n$ to a multiple of
$k$, and pad the coefficient columns to $k$. There are
$\lceil n/k\rceil^2$ square block products. The scalar operation
count is
\[
 O\!\left(n^2C^{\omega-2}+C^\omega\right).
\]
This estimate also covers $C>n$. Recursive input sums have at most
$O(\log C)$ levels, and output recombination adds at most four
terms per level. Even the crude bound of $C^4$ on intermediate
magnitude growth gives bit length
$M+b+O(\log n+\log C)=O_r(B+D)$. Schoolbook scalar arithmetic
therefore costs $O_r(M^2)$ per operation. The first affine layer
in basis coordinates has only $r$ inputs and is computed directly;
its rational coefficients introduce no dense matrix factor.

Nonlinear coefficient recurrences use $O_r(K^{2r})$ scalar
operations per neuron and center. Constant exponentials use
$O_r(M^3)$ bits. The full preprocessing bound is consequently
\[
 O_r\!\left(n^2B^3+D\left[
 n^2C^{\omega-2}+C^\omega+
 nD^{r/2}K^{2r}+nD^{r/2}M\right]M^2\right),
\]
where the first term is the basis computation.
Working matrices and stored tables fit this bound. Only the final
Taylor table, the network weights, and the scalar source selectors
are retained online.

At $r=3$ and logarithmic depth and parameter precision,
$C=O(\log^{9/2}n)$. The dense term is
\[
 O\!\left(n^2\log^{3+(9/2)(\omega-2)}n\right).
\]
Since $7^6<2^{17}$, one has $\omega<17/6$, and its logarithmic
exponent is less than $27/4<7$. The nonlinear term is
$O(n\log^{21/2}n)$ and the padding term is a fixed power of
$\log n$. Both fit $O(n^2\log^7n)$ because every fixed power
of $\log n$ is dominated by any fixed positive power of $n$.
This is $O(Dn^2\log^6n)$ for the stated depth. The rank-one mean
chain supplies the matching query lower bound.
\end{proof}

\begin{corollary}[The finite-input critical law for low input rank]
\label{cor:rank-three-finite-input}
For first-layer rank at most $r\in\{1,2,3\}$, the samplers can be
chosen to use
\[
 O\!\left(\min\{\sqrt D,n/\sqrt D\}\right)
\]
expected input probes and
\[
 O\!\left(
 (1+\min\{\sqrt D,n/\sqrt D\})
 B^{\max\{4,r+2\}}\right)
\]
expected bit operations, after the stated general-depth preprocessing.
The worst-case query order over this network class is sharp when
$b\ge\lceil\log_2\min\{n,D+1\}\rceil$.
\end{corollary}

\paragraph{Proof idea.}
The same function oracle also supports a rare full-input read. On the open
amplitude gate, reading the context determines the few basis means exactly.
The cached oracle then samples the known output probability without a dense
forward pass except on a sufficiently rare precision-refinement event.

\begin{proof}
In addition to the source factory, consider opening its dyadic gate
$q=O(D^{-1/2})$ and reading all $n$ signs only when that gate opens.
Compute the $r$ basis means exactly with $O(nB^2)$ bit operations.
The known sign mean on the open branch is $F(z)/q$. Its cached
interval costs $O(B^{\max\{4,r+2\}})$ bits. Direct refinement
has unconditional conditional-on-gate cost
$O(H\varepsilon B^4)\le O(B^4)$, because
$\varepsilon\le H^{-2}$ and $H=n^2D$. This follows from the same
fresh-uniform interval-measure argument as before. The expected
probe count is at most $nq$ and the expected bit cost is
$O((1+n/\sqrt D)B^{\max\{4,r+2\}})$. Choose this algorithm
when $n\le D$ and the source factory otherwise.

For the lower bound choose the largest power of two
$m\le\min\{n,D+1\}$, average those $m$ coordinates in the
first layer, and use the scalar critical mean chain. Its finite-input
lower bound is
$\Omega(\min\{\sqrt D,m/\sqrt D\})$, which is within an absolute
factor of the claimed expression. The remaining inputs and neurons
are ignored. This first matrix has rank one, and $1/m$ belongs to
the stated dyadic grid.
\end{proof}
